\section{Introduction}\label{sec:intro}

Probabilistic circuits (PCs) represent a distribution on $\{0,1\}^M$ by an arithmetic circuit with nonnegative weights. Smoothness and decomposability are the structural conditions under which marginals can be computed in time linear in the circuit size~\cite{ChoiVergariVdB20}. Which tractable distributions fall outside this class? We study the uniform spanning tree. By Kirchhoff's matrix-tree theorem~\cite{Kirchhoff1847} its generating polynomial is a determinant, so it has a polynomial-size generating circuit with subtraction, in the sense of Zhang, Juba and Van den Broeck~\cite{ZhangJubaVdB21}. Exact representation of uniform spanning trees of the complete graph by decomposable circuits needs exponential size~\cite{MartensMedabalimi14,ZhangHoltzenVdB20}. Those results concern exact representation or vanishing error on the complete graph, where the number of variables is quadratic in the number of vertices.

We prove lower bounds at constant total-variation error, without assuming determinism, for 2-connected bounded-degree expanders, the complete graph and dense graphs with a linear spectral gap. The two arguments behind these bounds are dual cases of one argument about matroids.

\subsection{Main result}

For a graph $G$, the edge expansion is $h(G)=\min\{\abs{\partial S}/\abs S:\ 0<\abs S\le\abs{V(G)}/2\}$, where $\partial S$ is the set of edges with exactly one endpoint in $S$. A graph is 2-connected if it has at least three vertices and stays connected after deleting any one vertex. Circuit terms are defined in Section~\ref{sec:models}. The size of a circuit is the number of nodes plus the number of wires.

\begin{theorem}[main theorem]\label{thm:main}
Fix an integer $D\ge2$ and a real $h>0$. Let $G$ be a 2-connected simple graph with $N$ vertices, $M$ edges, maximum degree at most $D$ and edge expansion $h(G)\ge h$, and let $P$ be the uniform distribution on the spanning trees of $G$, viewed as a distribution on $\{0,1\}^{E(G)}$. Let $C$ be a smooth decomposable probabilistic circuit with nonnegative weights over the $M$ edge variables, whose leaves have arity at most $2M/3$, and let $Q_C$ be its normalized distribution. If $\TV(P,Q_C)\le\eps<1$, then
\[
  \abs C\ \ge\ (1-\eps)^2\,e^{c N},\qquad
  c=c(D,h)=\frac{h_1^{5D+4}}{(D+1)^{4(D+1)}\,e^{10(D+1)}},\qquad h_1=\min(h,1).
\]
If moreover the Laplacian of $G$ has second smallest eigenvalue $\lambda_2(L_G)\ge\lambda>1$, the same holds with
\[
  c=c_\lambda(D)=\frac{\lambda\,(D-1)\,(1-1/\lambda)^D}{3D^3(D+\lambda/2)(D^2+1)} .
\]
\end{theorem}

The leaf-arity condition cannot be dropped: a single leaf of arity $M$ can store $P$. The bound $2M/3$ comes from the circuit reduction (Lemma~\ref{lem:decomp}). Girth plays no role. Theorem~\ref{thm:main} applies, with probability tending to one, to random $d$-regular graphs for every fixed $d\ge3$, to the Ramanujan graphs of Lubotzky, Phillips and Sarnak~\cite{LPS88}, and to the prisms and squares of bounded-degree expanders (Corollary~\ref{cor:families}).

\begin{corollary}[separation]\label{cor:separation}
Fix $D$, $h>0$, $\eps<1$, and a family of graphs $G$ as in Theorem~\ref{thm:main}.
\begin{enumerate}[label=(\alph*)]
\item Every nonnegative smooth decomposable PC with leaves of arity at most $2M/3$ that approximates $P$ within total variation $\eps$ has size $2^{\Omega(N)}=2^{\Omega(M)}$. Some such circuit has size $2^{O(M)}$, so $2^{\Theta(M)}$ is the optimal size.
\item The generating polynomial $\sum_xP(x)\prod_{e:x_e=1}z_e$ has an arithmetic circuit with subtraction of size polynomial in $N$, and every marginal probability $P(x_S=a)$ can be computed in polynomial time (Proposition~\ref{prop:upper}).
\end{enumerate}
\end{corollary}

For fixed $D$ the constant $c(D,h)$ is polynomial in $h$, with $\log(1/c)=O(D\log((D+1)/\min(h,1)))$; the previous version of this paper had $c\le D^{-4D/h}$, exponentially small in $D/h$. For $D=3$ and $h=1/5$ the closed form gives $c\approx5\cdot10^{-41}$ and the optimized parameters of Proposition~\ref{prop:touch-param} about $5\cdot10^{-24}$, compared with $9\cdot10^{-44}$ before (Table~\ref{tab:constants}). The spectral constant needs no small parameter at all: for 6-regular Ramanujan graphs, $\lambda_2\ge6-2\sqrt5$ gives $c\ge8.0\cdot10^{-8}$, and $c_\lambda(D)\ge2.7\cdot10^{-9}$ for all $d$-regular Ramanujan graphs with $6\le d\le64$. For Ramanujan graphs $c_\lambda(D)$ decays only polynomially in $D$, like $D^{-4}$. The degree of the polynomial in $h$ in $c(D,h)$ grows with $D$, and for the star bits behind it this is unavoidable: on 2-connected hub graphs they certify at most $O(h^{D/2+1})$ for some balanced split (Corollary~\ref{cor:hub}). A second, local argument gives much better constants on two explicit families. The prism $H\square K_2$ consists of two copies of $H$ joined by a perfect matching between corresponding vertices. The square $H^2$ has the vertices of $H$, with an edge between any two vertices at distance one or two in $H$.

\begin{corollary}[prisms and squares]\label{cor:prism-square}
Fix an integer $\Delta\ge2$ and a real $h_0>0$. Let $H$ be a connected simple graph with at least three vertices, maximum degree at most $\Delta$ and $h(H)\ge h_0$, and let $G$ be $H\square K_2$ or $H^2$. Then $G$ is 2-connected, and the conclusion of Theorem~\ref{thm:main} holds with
\[
  c=\frac{\min(1,h_0)}{24\,\Delta(\Delta+1)^5(3\Delta+4)}\quad\text{for }G=H\square K_2,\qquad
  c=\frac{8h_0}{243\,\Delta^{9}(\Delta-1)(h_0+\Delta^2)}\quad\text{for }G=H^2 .
\]
\end{corollary}

These constants come from Theorem~\ref{thm:general}, which applies to every bounded-degree expander whose vertex neighbourhoods are tied together by short cycles (Definition~\ref{def:linked}). For a cubic $H$ with $h(H)\ge1/5$, the prism gives $c\ge3\cdot10^{-7}$ (Example~\ref{ex:cubic}). The same local argument gives $e^{\Omega(L)}$ for the $L\times L$ grid (Corollary~\ref{cor:grid}), which is not an expander and is not covered by Theorem~\ref{thm:main}. All the constants are tiny, and in exact computations on small graphs the bounds we prove are far from tight (Appendix~\ref{app:checks}). The results are qualitative. For the bounded-degree expanders the exponent is linear in the number of variables; for the grid it is $\Omega(L)=\Omega(\sqrt M)$.

Expansion alone is not enough. The uniform spanning tree is the product of the uniform spanning trees of the blocks of $G$, so a balanced split along blocks has Bhattacharyya coefficient one with a product distribution. Proposition~\ref{prop:blocks} shows that such a split exists exactly when no block has more than $2M/3$ edges. Pendant edges give bounded-degree expanders of this kind (Proposition~\ref{prop:pendant}). For bounded-degree expanders this is the only obstruction: blocks inherit expansion, so if the largest block $G_1$ has $M_1>2M/3$ edges, every balanced split has $\beta^2\le e^{-\Omega(M_1-2M/3)}$, and for circuits with univariate leaves the bound $e^{\Omega(\abs{V(G_1)})}$ holds without any threshold (Theorem~\ref{thm:blocks}). 2-connectivity is the case $G_1=G$.

We also record exact-support consequences (Corollary~\ref{cor:support}). For every graph as in Theorem~\ref{thm:main}, every nonnegative smooth decomposable PC with the same leaf-arity condition, and every monotone arithmetic circuit, whose support is exactly the set of spanning trees of $G$, with arbitrary positive values on that set, has size $e^{\Omega(N)}$.

\subsection{The complete graph}

Both arguments behind Theorem~\ref{thm:main} need bounded degree. Let $\mu_n$ be the uniform spanning tree of $K_n$, a distribution on $M=\binom n2$ variables. Section~\ref{sec:complete} proves that for $n\ge3$ every balanced split $A|B$ of $E(K_n)$ satisfies
\[
  -\ln\beta(\mu_n;A|B)^2\ \ge\ \frac n{146}-\frac12\ln n-3
\]
(Theorem~\ref{thm:complete}), with an explicit sum $\Phi_n$ behind it whose exponential rate is $1/144.86$ per vertex. Consequently every smooth decomposable PC with nonnegative weights and leaves of arity at most $2M/3$ with $1-\TV(\mu_n,Q_C)\ge\delta>0$ has size at least $\delta^2e^{n/146}/(e^3\sqrt n)$ (Corollary~\ref{cor:complete}). The arity condition cannot simply be dropped, since a single leaf of arity $M$ represents $\mu_n$. A subset recursion gives an exact deterministic circuit of size $O(n^43^n)$ (Proposition~\ref{prop:complete-upper}), so at every fixed error the optimal size is $2^{\Theta(n)}=2^{\Theta(\sqrt M)}$ (Corollary~\ref{cor:complete-order}). Martens and Medabalimi~\cite{MartensMedabalimi14} proved a lower bound of the same order for exact representation and for approximation arbitrarily well; ours holds at every fixed total-variation error below one. The proved rate is still far from the truth. Exact computation for $6\le n\le9$ gives $-\ln\beta^2$ between $0.12n$ and $0.30n$ on the splits we tried, and a Monte Carlo estimate, which is an upper bound on the true rate and is not certified, gives about $0.17n$ for large $n$ (Remark~\ref{rem:kn-ceiling}).

The same method gives $2^{\Omega(n)}$ for every graph whose Laplacian has $\lambda_2\ge\kappa n$, with a rate of order $\kappa^4e^{-1.14/\kappa}$ (Theorem~\ref{thm:dense} and Corollary~\ref{cor:dense}). For $G(n,\frac12)$ this applies conditionally on the standard concentration of $\lambda_2$, which we do not re-derive. A minimum degree of $n/2$ is not enough for the method: two cliques joined by a perfect matching have a balanced split with $\beta^2>e^{-2}$ (Proposition~\ref{prop:two-cliques}).

\subsection{Matroids}

A spanning tree contains no cycle and meets every cut. The local argument uses the first fact and the vertex-touch argument the second. Section~\ref{sec:matroid} shows that the two arguments are the circuit and cocircuit cases of one argument for the uniform distribution on the bases of a matroid. Each element of a circuit is certified by a cut that meets the circuit in exactly two elements, and each element of a cocircuit by a cycle; an exchange argument turns such a frame into a conditional lower bound (Lemma~\ref{lem:mat-frame}). This gives a bound for every matroid in terms of compatible framed circuits and cocircuits (Theorems~\ref{thm:mat-local} and~\ref{thm:mat-pc}), and the bound is invariant under matroid duality (Proposition~\ref{prop:mat-dual}). For balanced matroids, which include the regular ones, a global version replaces frames by leverages (Theorem~\ref{thm:mat-global}). For a totally unimodular representation the leverages are effective resistances (Lemma~\ref{lem:mat-resist}), and applied to vertex stars this recovers Theorem~\ref{thm:touch-cut}. As an example outside graphs and their duals we obtain $2^{\Omega(n)}$ bounds for binary non-regular matroids obtained from linked expanders by attaching a Fano plane to every edge (Proposition~\ref{prop:fano}). The key lemma of the global version fails for the non-balanced binary matroid $S_8$ (Section~\ref{sec:limits}); whether the global version holds for all matroids is open.

\subsection{Subtraction}

The uniform spanning tree has a small circuit with subtraction. Section~\ref{sec:signed} asks how much of the lower bound survives when subtraction is allowed but the circuit keeps a fixed decomposition structure. For structured-decomposable circuits with real weights, a rank argument adapted from de Colnet and Mengel~\cite{DeColnetMengel21} gives exponential size for exact representation on linked expanders (Theorem~\ref{thm:signed-rank} and Corollary~\ref{cor:signed-exact}). One inequality for disjoint permutation blocks of $\UST$ (Theorem~\ref{thm:block-approx}) extends this to approximation. If the output $Q$ sums to $1$, every fixed error $\frac12\norm{\UST-Q}_1\le\eps<1$ needs size $e^{\Omega(N)}$, although $Q$ may be negative (Theorem~\ref{thm:signed-o2} and Corollary~\ref{cor:signed-approx}); the key step is that a matrix of rank $r$ and bounded $\ell_1$ mass covers at most a $\sqrt{r/n}$ fraction of an $n\times n$ permutation, a consequence of a rank bound of Ky Fan and Hoffman~\cite{FanHoffman53}. Without normalization the threshold is $\eps<\frac12$, which is sharp because $Q=0$ has error $\frac12$ (Theorem~\ref{thm:signed-approx}). The same bounds hold on every 2-connected bounded-degree expander, at any girth and without smoothness (Corollary~\ref{cor:signed-expander}), with blocks built from pendant vertices instead of short cycles, and they cover sparse and low-degree approximations (Corollary~\ref{cor:signed-sparse}). Normalized squares and positive semidefinite forms of structured-decomposable amplitudes obey the same lower bound as nonnegative circuits (Theorem~\ref{thm:signed-square}): for structured circuits the barrier is decomposability, not sign. For unstructured decomposable circuits with subtraction we have lower bounds only for restricted classes. On bounded-degree expanders, support-exact representation needs size $2^{N^{\Omega(1/d)}}$ at product-depth $d$ and $N^{\Omega(\log N)}$ for formulas, and sums of $r$ structured circuits, each on its own vtree, need total size $2^{\Omega(3^{-r}\sqrt N)}$ (Corollary~\ref{cor:signed-unstructured}). These proofs transfer the restriction arguments of Raz~\cite{Raz09} and of Raz and Yehudayoff~\cite{RazYehudayoff09} through a large complete minor~\cite{KrivelevichNenadov21,Krivelevich19}, after checking that their rank bounds do not depend on the constants of the restriction. For sums of $r$ structured circuits, a packing of far-apart pendant gadgets gives total size $2^{\Omega_r(N)}$ directly in the graph, both for support-exact representation and at error $\eps<\frac12$ (Theorem~\ref{thm:signed-sums-linear}); its constants are astronomically small, so the minor bound remains better at any realistic size. General circuits stay open. For exact representation the general question is equivalent to the syntactically multilinear complexity of the spanning-tree polynomial (Proposition~\ref{prop:signed-multilinear}), and a small approximating circuit would have to cancel all but an exponentially small fraction of its absolute mass (Theorem~\ref{thm:signed-cancel}), largely between terms on different splits (Proposition~\ref{prop:signed-cross}). We explain why the block method does not yet reach unstructured circuits (Section~\ref{sec:signed-barriers}).

\subsection{What is known and what is new}\label{sec:intro-priority}

The ingredients of the proof are known; we credit them where they are used.
\begin{itemize}
\item For a joint distribution $\mu$ of $(X,Y)$, the largest Bhattacharyya coefficient between $\mu$ and a product distribution equals the top singular value of the entrywise square root of $\mu$. This is Lemma~6 of Lapidoth and Pfister~\cite{LapidothPfister19}, stated for their order-$\frac12$ dependence measure $J_{1/2}$.
\item The adaptive chain rule we use (Lemma~\ref{lem:chain}) is a case of Exercise~I.42(b) in Polyanskiy and Wu~\cite{PolyanskiyWu25}, combined with the previous item. For independent pairs it is the additivity of $J_{1/2}$~\cite[Lemma~12]{LapidothPfister19}. Hellinger distance and its product structure are standard tools for communication lower bounds~\cite{BJKS04,Jayram09}.
\item The reduction from circuits to products across balanced edge splits is the balanced decomposition and the first step of the overlap lemma from our earlier paper on dimers~\cite[Lemmas~7.1 and~7.2]{Blaschke26Dimers}. The decomposition itself is the standard balancing argument for syntactically multilinear circuits~\cite{RazYehudayoff08}, also stated for sum-product networks by Martens and Medabalimi~\cite[Thm.~38]{MartensMedabalimi14}.
\item Kirchhoff's formula for edge marginals, Rayleigh monotonicity, Thomson's principle and the triangle inequality for effective resistance are classical~\cite{Kirchhoff1847,LyonsPeres16,KleinRandic93}, as is Mohar's Cheeger-type inequality~\cite{Mohar89}, whose proof we repeat for the lower bound we use (Lemma~\ref{lem:cheeger}). The count of connected sets (Lemma~\ref{lem:count}) is a form of a lemma of Fomin and Villanger~\cite{FominVillanger12}, and the random deletion (Lemma~\ref{lem:delete}) is the alteration method with a Peierls-type union bound. Lemma~\ref{lem:forest-mgf} and Section~\ref{sec:edge-pairs} use the transfer-current theorem of Burton and Pemantle~\cite{BurtonPemantle93}. The proof of Theorem~\ref{thm:complete} uses de Caen's bound on sums of squared degrees~\cite{deCaen98}, the exact law of the isolated vertices of a uniform labelled forest (a Cayley-type count), and Robbins' form of Stirling's formula~\cite{Robbins55}.
\item For matroids we use standard facts from Oxley~\cite{Oxley11}, the negative correlation of regular matroids due to Feder and Mihail~\cite{FederMihail92}, and the characterization of balanced binary matroids by Choe and Wagner~\cite{ChoeWagner03}.
\end{itemize}
Constant-error lower bounds against nonnegative decomposable circuits are not new in general. Leland and Choi~\cite{LelandChoi25} prove one for \emph{deterministic} circuits. A generic $\exp(\Omega(n))$ separation from a random linear code, whose inference is easy by Gaussian elimination, follows from the rank bounds of de Colnet and Mengel~\cite{DeColnetMengel20}, as explained in~\cite[Sec.~10.4]{Blaschke26Dimers}. The dimer paper~\cite{Blaschke26Dimers} proves an $\exp(\Theta(\sqrt n))$ bound for the uniform domino tiling of the square grid. That paper also lists the robust spanning-tree question for expanders as open~\cite[Sec.~11, problem~5]{Blaschke26Dimers}.

We believe the following are new.
\begin{enumerate}[label=(\roman*)]
\item The constant-error $2^{\Omega(N)}$ lower bound against nonnegative smooth decomposable PCs, without determinism, for uniform spanning trees of every 2-connected bounded-degree expander, at any girth. Apart from an earlier version of this paper, which covered prisms and squares of expanders, we know of no constant-error exponential lower bound against nondeterministic decomposable PCs for uniform spanning trees of any family, and none for a distribution of this kind whose inference reduces to determinants.
\item The vertex-touch reduction: two bits per vertex that has edges on both sides of the split, conditioning only on the stars of a sparse set of such vertices, and a conditional lower bound from effective resistances, in a telescoped form controlled by $\lambda_2$ of the graph left after the sparse set is deleted, with a random deletion that keeps this graph an expander (Lemma~\ref{lem:touch}, Theorem~\ref{thm:touch-cut}, Lemmas~\ref{lem:touch-flow} and~\ref{lem:delete}).
\item The local cycle-bit reduction: for spanning trees, two indicator bits per crossing cycle, together with the spatial Markov property at the edges incident to the cycle, give the conditional non-degeneracy that the chain rule needs (Lemma~\ref{lem:local} and Theorem~\ref{thm:cut}). It gives the better constants of Corollary~\ref{cor:prism-square} and the grid bound.
\item The constant-error $2^{\Omega(n)}$ bound for the uniform spanning tree of $K_n$ with an explicit rate (Theorem~\ref{thm:complete} and Corollary~\ref{cor:complete}), its extension to dense graphs (Corollary~\ref{cor:dense}), and the averaged conditioning on a rectangle (Lemma~\ref{lem:avg-cond}) that lets the proof use a conditioning set that is not independent. We did not find Lemma~\ref{lem:avg-cond} stated in this form, but it is elementary.
\item The matroid form of both arguments: framed circuits and cocircuits, the local theorem for every matroid, the global theorem for balanced matroids, and the Fano-decorated family (Section~\ref{sec:matroid}). We did not find the exchange bound of Lemma~\ref{lem:mat-exchange} stated for matroids, but it is also elementary.
\item The approximate lower bounds for structured-decomposable circuits with real weights: the permutation-block inequality (Theorem~\ref{thm:block-approx}), which reaches every error below $1$ for normalized outputs, and its pendant-vertex form on every 2-connected bounded-degree expander, without smoothness (Corollary~\ref{cor:signed-expander}). Lemma~\ref{lem:kfh} is classical~\cite{FanHoffman53}.
\item The support-exact lower bounds for restricted unstructured signed circuits (Corollary~\ref{cor:signed-unstructured}), and the bounds for sums of structured circuits through packings of pendant gadgets (Theorem~\ref{thm:signed-sums-linear}). We did not find them in the literature, but the transfer is short and may be folklore; the rank arguments are those of~\cite{Raz09,RazYehudayoff09}.
\item The limits of the vertex-touch method: on 2-connected hub graphs with expansion $h$ and maximum degree $D$, star bits certify at most $O(h^{D/2+1})$ for one balanced split, while bits for pairs of edges at a vertex certify $h^3/(10935D^5)$ for that split (Corollary~\ref{cor:hub} and Proposition~\ref{prop:hub-pairs}). No novelty is claimed for the ingredients.
\end{enumerate}
Our literature search was targeted, not exhaustive. The closest results we found are discussed in Section~\ref{sec:related}. In particular, Chattopadhyay, Datta and Mukhopadhyay~\cite{CDM20Negations} and Chattopadhyay, Datta, Ghosal and Mukhopadhyay~\cite{CDGM22} prove exact $2^{\Omega(n)}$ monotone lower bounds for directed spanning-tree polynomials of expanders of large constant degree, and M\'esz\'aros and Peled~\cite{MeszarosPeled26} prove that the uniform spanning tree of a graph with large minimum degree is close to a product-like process in expected edit distance. Neither statement concerns total variation for circuits; Section~\ref{sec:related} explains why neither contradicts nor implies ours.

\subsection{Proof outline}\label{sec:outline}

Fix a split $E=A\sqcup B$ of the edges. For a product distribution $q=q^A\otimes q^B$ write $\BC(P,q)=\sum_x\sqrt{P(x)q(x)}$, and let $\beta(P;A|B)$ be the largest value over all products. The circuit reduction (Section~\ref{sec:circuits}) gives
\[
  \abs C\ \ge\ (1-\eps)^2\Big/\max_{A|B\ \text{balanced}}\beta(P;A|B)^2 .
\]
So it suffices to show $\beta\le e^{-cN/2}$ for every split in which both sides have at least $M/3$ edges. The proof of Theorem~\ref{thm:main} has four steps.
\begin{enumerate}
\item \emph{Vertex-touch bits.} Call a vertex $w$ mixed if it has edges in both $A$ and $B$. Put $a_w=\one[T\text{ uses no }A\text{-edge at }w]$ and $b_w=\one[T\text{ uses no }B\text{-edge at }w]$. The bits $a$ depend only on $T\cap A$ and the bits $b$ only on $T\cap B$. A spanning tree has an edge at every vertex, so $a_w=b_w=1$ never happens. Data processing gives $\beta(P;A|B)\le\opnorm{\sqrt\mu}$, where $\mu$ is the joint law of the bits at a set $W$ of mixed vertices (Lemma~\ref{lem:dpi}).
\item \emph{Chain rule.} If, given every possible value of the earlier pairs, the pair $(a_w,b_w)$ has probability at least $p_A$ of being $(1,0)$ and at least $p_B$ of being $(0,1)$, then the pair contributes a factor $((1+\sqrt{1-4p_Ap_B})/2)^{1/2}\le e^{-p_Ap_B/2}$ to $\opnorm{\sqrt\mu}$ (Lemmas~\ref{lem:chain} and~\ref{lem:two-by-two}). No independence between the pairs is needed.
\item \emph{Conditioning on stars.} Take $W$ independent and condition on the tree on the stars of the other vertices of $W$. The rest of the tree is a uniform spanning tree of a graph obtained from $G-(W\setminus\{w\})$ by identifying vertices. Deleting the $A$-edges at $w$ one at a time, Kirchhoff's formula and Rayleigh monotonicity give $\Pr[a_w=1\,|\,\text{past}]\ge\prod_u(1+R_u)^{-1}$, with effective resistances $R_u$ in a graph that contains $G-W$ (Lemma~\ref{lem:touch}). Spreading the current from $u$ over all remaining neighbours of $w$, the product telescopes to $\Pr[a_w=1\,|\,\text{past}]\ge(b/d)(1+1/\lambda_2(L_{G-W}))^{-a}$, where $a$ and $b$ are the numbers of $A$- and $B$-edges at $w$ and $d=a+b$ (Lemma~\ref{lem:touch-flow}).
\item \emph{Expansion.} Every balanced split of an expander has linearly many mixed vertices (Lemma~\ref{lem:mixed-vertices}). If $G$ is 2-connected, a random sparse set $W$ of them, with the vertices of a few dangerous configurations removed, leaves $G-W$ an expander (Lemma~\ref{lem:delete}); the dangerous configurations are counted by a Peierls argument over connected sets (Lemma~\ref{lem:count}). Mohar's inequality turns the expansion of $G-W$ into a bound on $\lambda_2$ (Lemma~\ref{lem:cheeger}). If $\lambda_2(L_G)>1$, a set $W$ with pairwise distances at least $3$ suffices, since deleting it lowers $\lambda_2$ by at most one (Lemma~\ref{lem:sep3}).
\end{enumerate}
The local argument for Corollary~\ref{cor:prism-square} replaces step~1 by two bits per short cycle that crosses the split ($u=\one[\text{all }A\text{-edges of the cycle are in }T]$ and $v$ likewise for $B$; a tree contains no cycle), step~3 by conditioning on everything outside the edges incident to the cycle and counting completions (Lemma~\ref{lem:local}), and step~4 by a packing of short cycles (Lemma~\ref{lem:mixed}). Cycles and cuts give the two arguments: a tree contains no cycle and meets every cut. Section~\ref{sec:matroid} turns this remark into a theorem about matroids. Only steps~1 and~3 use spanning trees.

For $K_n$ (Section~\ref{sec:complete}) the set $W$ cannot be independent. We condition on the tree edges inside $W$, keep the set $I$ of vertices of $W$ that these edges do not touch, and average the resulting bounds over the conditioning (Lemma~\ref{lem:avg-cond}). Given the conditioning, a star-deletion bound that is sharp on complete graphs gives the kernel of each vertex of $I$ (Lemma~\ref{lem:kn-star}). The law of $I$ for a uniform $r$-set $W$ is explicit (Lemma~\ref{lem:kn-iso}), de Caen's inequality shows that an average vertex has a constant fraction of its edges on each side, and choosing $r\approx0.38n$ gives the rate. For dense graphs, the transfer-current theorem and Hadamard's inequality replace the exact law of $I$ (Lemma~\ref{lem:forest-mgf}). The determinant enters only on the upper-bound side.

\paragraph{Organization.} Sections~\ref{sec:prelim} and~\ref{sec:tools} give notation, circuit models and the known tools. Section~\ref{sec:vertex} proves the bound for 2-connected expanders, and Section~\ref{sec:complete} adapts it to the complete graph. Sections~\ref{sec:trees} and~\ref{sec:geometry} give the local cycle-bit argument. Section~\ref{sec:main} proves the main results, with constants, the families covered, the upper bound and exact-support consequences. Section~\ref{sec:matroid} puts both arguments in the setting of matroids. Section~\ref{sec:signed} treats circuits with subtraction. Section~\ref{sec:limits} discusses limitations and variants, including the block structure and the matroid $S_8$; Sections~\ref{sec:related} and~\ref{sec:open} cover related work and open questions. Appendix~\ref{app:checks} describes the machine checks.
